\begin{theorem}
	If \( F \) is a finitely generated \underline{free} module over \( R \) and \( G \subseteq F \) a
	submodule \( \implies \exists \) basis \( w_1,\ldots, w_n \) of \( F \), and \( a_1,\ldots,a_m \in
	R\) nonzero, \( m\leq n \), \( a_1 \mid a_2 \mid \cdots \mid a_m \), such that \( a_1w_1,\ldots,
	a_m w_m \) is a basis of \( G \). In particular, \( G \) is free.\todo{Proof gives later.}
\end{theorem}

\begin{proof}
	\( F \) is finitely generated over \( R \) with \( R \) being PID and thus Noetherian, then \( F \) is also
	a Noetherian \( R \)-module, thus \( G \) is finitely generated. In particular, we may consider
	\[
		\rank(G) \coloneqq \dim_{S^{-1}R}S^{-1}G < \infty
	\]
	where \( S= R \backslash \{0\} \). We shall argue by induction on \( m = \rank(G) \) to show the
	following:
	\begin{itemize}
		\item \( G \) is free.
		\item Then the full statement of the theorem holds.
	\end{itemize}
	\begin{enumerate}
		\item \( m = 0 \), this means that \( \forall \; u \in G, \frac{u}{1} = 0 \) in \( S^{-1} G \),
		      \( \iff \; \exists \; a \in R \backslash \{0\} \), s.t. \( au = 0 \). Fix a basis \(
		      u_1,\ldots,u_n \) of \( F \), if we write
		      \[
			      u = \alpha_1 u_1 + \cdots + \alpha_n u_n
		      \]
		      with \( au = 0 \implies \)
		      \[
			      \sum_{i}(a\alpha_i) u_i = 0 \implies a\alpha_i = 0 \; \forall \; i
		      \]
		      since \( a \ne 0 \), thus by domain we have \( \alpha_i = 0\; \forall \; i \implies u=0
		      \). Hence \( G = \{0\} \), and the conclusion is clear.
		\item Inductive step: Assumpe that \( m > 0 \implies G \ne \{0\} \). Consider the \( R
		      \)-linear maps \( \vp : F \to R \), note that this is nonempty as we can take the \( 0 \)
		      map. Note that \( \vp(G) \) is ideal of \( R \). Since \( R \) is Noetherian, there exists
		      \( \vp \) such that \( \vp(G) \) is maximal among all such choices of \( \vp \) w.r.t.
		      inclusion. \( R \) is a PID implies that \( \vp(G) = (a_i) \) for some \( a_i \in R \).
		      Can choose \( x_1 \in G \), s.t. \( \vp(x_1) = a_1 \).
		      \begin{claim}
			      \( a_1 \ne 0 \)
		      \end{claim}
		      \begin{proof}
			      A linear map \( \psi : F \to R \) is uniquely determined by giving \( \psi (u_i), 1 \leq
			      i \leq n\). In particular, have
			      \[
				      \begin{aligned}
					      \pi_j : F  & \to R \\
					      \pi_j(u_i) & =
					      \begin{cases}
						      1, \text{ if } i  = j \\
						      0, \text{ otherwise.}
					      \end{cases}
				      \end{aligned}
			      \]
			      if \( a_1 = 0 \), by maximality of \( a_1 \implies \psi(G) = \{0\} \; \forall \; \psi :
			      F \to R\) being \( R \)-linear map. Thus \( \pi_j(G) = \{0\} \; \forall \; j \implies G
			      = \{0\}\), leading to contradiction with our inductive asssumption \( \lightning \).
			      Thus \( a_1 \ne 0 \).
		      \end{proof}
		      \begin{claim}
			      If \( \psi: F \to R \) arbitrarily linear map, then \( a_1 \mid \psi(x) \) where \(
			      \psi(x) = a_1, x \in G \).
		      \end{claim}
		      \begin{proof}
			      Let \( b = \gcd\{a_1,\psi(x)\}  \xRightarrow{ R \text{ PID}} b = \alpha \underbrace{a_1}_{\vp(x)} + \beta
			      \psi(x)\) for some \( \alpha, \beta \in R \). Now if \( \vp' = \alpha \vp + \beta \psi
			      \implies \vp'(x) = b\) and thus \( b \mid a_1 \implies (\vp'(x)) \supseteq (\vp(x))
			      \). By maximality: \( (\vp'(x))  = (\vp(x)) = (b) \implies b \sim a_1\) but \( b \mid
			      \psi(x) \implies a_1 \mid \psi(x)\).
		      \end{proof}
		      Now apply this for \( \psi = \pi_j \) for each \( j \), we have
		      \[
			      \pi_j (x) = a_1 c_j \text{ for some } c_j \in R \; \forall \; j
		      \]
		      and thus
		      \[
			      x = \sum_{j=1}^n \pi_j (x) u_j = a_1 \cdot \underbrace{\sum_{j} c_j u_j}_{e_1}
			      \implies x = a_1 e_1 \text{ for some } e_1 \in F
		      \]
		      and thus
		      \[
			      a_1 = \vp(x) = \vp(a_1 e_1) = a_1 \vp(e_1)
		      \]
		      by that \( a\ne 0 \) we have \( \vp(e_1) = 1 \).

		      \begin{claim}
			      Let \( K = \ker(\vp) \), we have \( F = K \oplus R e_1 \).
		      \end{claim}

		      \begin{proof}
			      \leavevmode
			      \begin{itemize}
				      \item \( K \cap R e_1 = \{0\} \): if \( \lambda e_1 \in K \implies \)
				            \[
					            \lambda = \lambda \vp(e_1) = \vp(\lambda e_1) = 0 \implies \lambda e_1 = 0
				            \]
				      \item \( K + Re_1 = F \): if \( u \in F \implies u - \vp(u)e_1 \in K \) since \(
				            \vp(u - \vp(u) e_1) = \vp(u) - \vp(u)\vp(e_1) = 0 \implies F = K + Re_1 \).
			      \end{itemize}
		      \end{proof}

		      \begin{claim}
			      \( G = (K \cap G) \oplus R a_1 e_1 \).
			      \begin{note}
				      \( a_1 e_1 = x \in G \).
			      \end{note}
		      \end{claim}

		      \begin{itemize}
			      \item \( (K \cap G) \cap R a_1 e_1 = \{0\} \) is clear by previous reasoning, basically
			            larger thing already intersect with only \( 0 \).
			      \item \( G = (K \cap G) + R a_1 e_1 \): have \( \vp(G) = (a_1) \), if \( u \in G \implies  \)
			            \[
				            u - \underbrace{\underbrace{\vp(u)}_{(a_1)}e_1}_{\in R a_1 e_1} \in K \cap G
			            \]
			            since \( u \in G, a_1 e_1 \in G, u - \vp(e_1) \in K \) as above.
		      \end{itemize}

		      \begin{conclusion}
			      We have
			      \[
				      \begin{aligned}
					      F & = K \oplus R e_1              \\
					      G & = (K \cap G) \oplus R a_1 e_1
				      \end{aligned}
			      \]
		      \end{conclusion}
		      We have that \( a_1 \ne 0, x = a_1 e_1 \implies e_1 \ne 0 \implies \rank(Ra_1 e_1) \ne
		      0\) and actually \( \rank(R a_1 e_1 ) = 1 \). Thus we have
		      \[
			      \rank(K \cap G) = \rank(G) - 1 = m - 1
		      \]
		      Apply induction on the assertion that ``every submodule of a free finitely generated
		      \( R \)-module is free'' and have \( K \cap G \) is free of rank \( m - 1 \). By
		      combining a basis of \( K \cap G \) with \( a_1 e_1 \implies \) get a basis of \( G
		      \). Hence we know by induction that all submodules of a free finitely generated \( R \)-module
		      are free. Now let's do induction for the full statement of the thream. Since \( K
		      \subseteq F \) is submodule, we know \( K \) is free. Use induction for \( K \cap G
		      \subseteq K \implies \exists \) basis \( e_2, \ldots, e_n \) of \( K \) and \(
		      a_2,\ldots, a_m \in R \) non-zero, such that
		      \[
			      \begin{cases}
				      a_2 e_2,\ldots, a_m e_m \text{ is a basis of } K \cap G \\
				      a_2 \mid a_3 \mid \cdots \mid a_m
			      \end{cases}
		      \]
		      Now \( F = K \oplus Re_1 \implies e_1,\ldots, e_n \) basis of \( F \). \( G = (K \cap
		      G)  \oplus R a_1 e_1 \implies a_1 e_1,\ldots, a_m e_m\) basis of \( G \). It is only
		      left to show:
		      \[
			      \text{If } m\geq 2 \implies a_1 \mid a_2
		      \]
		      Let \( d = \gcd(a_1,a_2) \xRightarrow{R \text{ PID}} d = \gamma a_1 + \delta a_2\). If \(
		      \psi: F \to R\) is s.t.
		      \[
			      \begin{cases}
				      \psi(e_1) = \gamma \\
				      \psi(e_2) = \delta \\
				      \psi(e_i) = 0 \text{ for } i > 2
			      \end{cases}
		      \]
		      Have \( \psi(a_1 e_1 + a_2 e_2) = d \implies \psi(G) \supseteq (d)\supseteq (a_1) \) and
		      since \( d \mid a_1 \implies \) by maximality of \( (a_1) \) we have \( (d) = (a_1)
		      \implies a_1 \mid a_2 \).
	\end{enumerate}
\end{proof}
The drawback of this theorem is that it doesn't give us a way to compute the \( a_1 e_1,\ldots,
a_m e_m \), but it is
good enough for our purpose. A more general theorem of this version will be proved in homework.

